\subsection{Ako spúšťam proces}
Celý pipeline spracovania dát sa skladá z nasledujúcich krokov, ktoré vykonávam postupne:

\subsubsection*{Krok 1 — Preklad nanopórových dát pomocou Dorado}
Surové nanopórové dáta preložím pomocou basecallera Dorado do formátu \texttt{.bam}:
\begin{verbatim}
dorado basecaller <model> <pod5_vstup> > vystup.bam
\end{verbatim}

\subsubsection*{Krok 2 — Konverzia BAM súboru do textového formátu}
Spustím skript \texttt{CloverGenerateFyle.py}, ktorý BAM súbor preloží do textového formátu \texttt{.txt}, ktorý následne môžem použiť ako vstup pre Clover:
\begin{verbatim}
python CloverGenerateFyle.py --input vystup.bam --output data.txt
\end{verbatim}

\subsubsection*{Krok 3 — Nájdenie primeru}
Zavolám skript \texttt{GettingPrimer.py}, ktorý na textový súbor zavolá Clover, nájde najväčší klaster a ten predá môjmu multi-klastrovaciemu algoritmu. Z neho vytiahnem najväčší klaster a na ten zavolám vlastný multi-alignment algoritmus. Výsledkom je konsenzuálna sekvencia primeru, ktorá sa uloží do súboru \texttt{primerFound.txt}:
\begin{verbatim}
python GettingPrimer.py --input data.txt --output primerFound.txt
\end{verbatim}

\subsubsection*{Krok 4 — Rozdelenie sekvencií podľa primeru}
Spustím skript \texttt{splitDnaByPrimer.py}, ktorý vyhľadáva primer v každej sekvencii textového súboru. Nájdené primery sa ukladajú do adresára \texttt{PrimersFound}:
\begin{verbatim}
python splitDnaByPrimer.py --input data.txt --primer primerFound.txt \
    --output PrimersFound/
\end{verbatim}

\subsubsection*{Krok 5 — Vyčistenie sekvencií od primeru}
Zavolám skript \texttt{SelectFromPrimer.py}, ktorý z predchádzajúceho textového výstupu vytvorí nový \texttt{.txt} súbor s DNA sekvenciami očistenými od primeru:
\begin{verbatim}
python SelectFromPrimer.py --input PrimersFound/ --output data_clean.txt
\end{verbatim}

\subsubsection*{Krok 6 — Záverečné klastrovanie a zarovnanie}
Nakoniec zavolám skript \texttt{postPrimerClover.py}, ktorý z dát očistených od primeru vytvorí Clover súbor, nájde klastre a začne zarovnávať. Pokiaľ nájde klastre požadovanej veľkosti (viac ako 10 nukleotidov), výsledky uloží do súborov vo formáte:
\begin{center}
\texttt{Aligned\_dnaSize-\{dist\}num-\{ClusterSize\}id-\{count\}.txt}
\end{center}
\begin{verbatim}
python postPrimerClover.py --input data_clean.txt --output Aligned_
\end{verbatim}

Celý proces je znázornený na diagrame nižšie — každý krok nadväzuje na výstup predchádzajúceho, pričom finálnym výsledkom sú zarovnané DNA reťazce roztriedené do klastrov.
